\documentclass[10pt,a4paper]{amsart}
\usepackage[margin=19mm]{geometry}
\usepackage{amsmath,amssymb,mathtools}
\usepackage[scr=rsfs,cal=euler]{mathalfa}
\usepackage{xfrac}
\usepackage[unicode,hidelinks,bookmarks=false]{hyperref}
\newcommand{\ie}{i.e.\ }
\newcommand{\cf}{cf.\ }
\newcommand{\resp}{respectively\ }
\newcommand{\Subsection}{\subsection}
\providecommand{\nobreakdash}{\nobreak\hbox{-}}
\setlength{\emergencystretch}{2em}
\multlinegap0pt
\usepackage{mathtools}
\newtheorem{lemma}{Lemma}
\newtheorem{theorem}{Theorem}
\title{Symmetric Bounded Indistinguishability: Hypergeometric Smoothing and Hahn Polynomials}
\author{Christopher Williamson}
\date{Source: arXiv:2605.13771v2 (CC BY 4.0)}
\begin{document}
\maketitle

\noindent\emph{Source and licence.} Symmetric Bounded Indistinguishability: Hypergeometric Smoothing and Hahn Polynomials. Version arXiv:2605.13771v2 (CC BY 4.0), distributed by arXiv under the Creative Commons Attribution 4.0 International licence, http://creativecommons.org/licenses/by/4.0/, as recorded in the arXiv licence metadata for this exact version and shown on the abstract page https://arxiv.org/abs/2605.13771v2. Full licence notice: \texttt{licenses/arxiv-2605.13771-CC-BY-4.0.txt}.\par\smallskip

\noindent\textit{Editorial scope.} Counted unit: the article's theorem thm:symmetric-repair-lower (source0.tex lines 1461--1481). The article's complete original proof of it is delivered with it (source0.tex lines 1750--1932, one \texttt{proof} environment of 183 lines, which the article places away from the statement). No mathematics is written, repaired or re-derived: the statement and the proof are the article's own lines, in the article's own notation, and no retained line is substituted or re-flowed.  Retained uncounted as the source-local context the proof needs: lines 663--668; lines 698--711; lines 1743--1746.  Nothing was omitted from any retained span, so the package carries no omission record.  No retained line contains a citation, so the wrapper carries no bibliography and no bibliography entry.  No retained line contains a figure environment, an image-inclusion command or drawing code, so no image is delivered and the package declares no asset dependency.  Editorial formatting only, in addition: an AMS \texttt{amsart} preamble that restates the article's own declarations and macros used by the retained lines, namely the environments \texttt{lemma}, \texttt{theorem} on separate counters, together with the standard packages those lines need.  The retained environments print in this document as Lemma 1, Theorem 1 in order of appearance, whereas the article prints the same items with the numbering of its own full document.  The complete native source of the article as distributed in the arXiv e-print for this exact version is bundled unchanged as \texttt{sources/200537/source0.tex}.
\begin{equation}
\label{eq:T-M-adjoint-measures}
    \sum_{a=0}^t f(a)(M_{n\to t}\sigma)(a)
    =
    \sum_{j=0}^n (T_{n,t}f)(j)\sigma(j).
\end{equation}

\begin{lemma}
\label{lem:hahn-measure-intertwining}
Let \(0\le t\le n\) and \(0\le r\le n\). Then
\[
    M_{n\to t}\theta_{n,r}
    =
    \begin{cases}
        \lambda_{n,t,r}\theta_{t,r},
            & 0\le r\le t,\\[1ex]
        0,
            & t<r\le n.
    \end{cases}
\]
\end{lemma}

\begin{lemma}
\label{lem:l1-ratio-simple-lower}
Let \(1\le k\le n\). Then $\frac{L_k}{L_n}\ge\frac{2^k}{\sqrt{(k+1)\binom{2k}{k}}}=\Theta(k^{-1/4})$.
\end{lemma}

\begin{theorem}[Repair lower bound]
\label{thm:symmetric-repair-lower}
Let \(1\le k<n\) and for some sufficiently small constant $c$, set
\[
    \delta_0
    :=
    \min\left\{
        1,\,
        \lambda_{n,k,k}\cdot\frac{c}{k^{1/4}}
    \right\}.
\]
Then for every \(0<\delta\le \delta_0\), there exist symmetric
distributions \(\mu,\nu\) on \(\{0,1\}^n\) satisfying $\Delta_k(\mu,\nu)=\delta$
and hence \((k,\delta)\)-wise indistinguishability, such that
\[
    \operatorname{Rep}_{n,k}(\mu,\nu)
    \ge
    \frac{\delta}
         {2\lambda_{n,k,k}\sqrt{n+1}}.
\]
\end{theorem}

\begin{proof}(of Theorem~\ref{thm:symmetric-repair-lower})
Define the signed measure
\[
    \theta:=\lambda_{n,k,k}^{-1}\theta_{n,k},
\]
from which we have
\[
    |\theta|(\{0,\dots,n\})=\lambda_{n,k,k}^{-1}L_n.
\]
Since \(k\ge 1\), the function \(\phi_k^{(n)}\) is orthogonal to the
constant function. Hence $\theta(\{0,\dots,n\})=0$.

Let
\[
    \gamma:=\frac{2\delta}{L_k}.
\]
By the assumption \(\delta\le\delta_0\) and Lemma~\ref{lem:l1-ratio-simple-lower}, we have
\[
    \gamma
    \le
    \frac{2\lambda_{n,k,k}}{L_n}
    =
    \frac{2}{|\theta|(\{0,\dots,n\})}.
\]

Define a probability measure \(\rho\) on \(\{0,\dots,n\}\) by
\[
    \rho(j):=\frac{|\theta(j)|}{|\theta|(\{0,\dots,n\})}.
\]
Now define probability measures \(\overline{\mu},\overline{\nu}\) on
\(\{0,\dots,n\}\) by
\[
    \overline{\mu}:=\rho+\frac{\gamma}{2}\theta,
    \qquad
    \overline{\nu}:=\rho-\frac{\gamma}{2}\theta.
\]
These measures have total mass \(1\), because \(\rho\) has total mass
\(1\) and \(\theta\) has total mass \(0\). They are nonnegative because,
for every \(j\),
\[
    \rho(j)-\frac{\gamma}{2}|\theta(j)|
    =
    |\theta(j)|
    \left(
        \frac{1}{|\theta|(\{0,\dots,n\})}
        -
        \frac{\gamma}{2}
    \right)
    \ge 0.
\]
Let \(\mu,\nu\) be the symmetric distributions on
\(\{0,1\}^n\) corresponding to $\overline{\mu}$ and $\overline{\nu}$.

Their signed Hamming-weight difference is
\[
    \sigma:=\overline{\mu}-\overline{\nu}
    =
    \gamma\theta.
\]
It follows immediately from the linearity of the marginalisation operator and Lemma~\ref{lem:hahn-measure-intertwining} that $M_{n\to k}\theta=\theta_{k,k}$ and therefore $M_{n\to k}\sigma=\gamma\theta_{k,k}$.

Hence
\[
    \Delta_k(\mu,\nu)=
    \frac12 |M_{n\to k}\sigma|(\{0,\dots,k\})=
    \frac{\gamma}{2}|\theta_{k,k}|(\{0,\dots,k\}) =
    \frac{\gamma}{2}L_k =
    \delta.
\]
Thus \(\mu,\nu\) are \((k,\delta)\)-wise indistinguishable.

We now show that any repair to obtain exact \(k\)-wise indistinguishability requires global adjustment of at least the claimed amount. Let \(\mu',\nu'\) be any symmetric
\((k,0)\)-wise indistinguishable distributions on \(\{0,1\}^n\), and let
\(\overline{\mu}',\overline{\nu}'\) be their Hamming-weight laws. Define
\[
    \sigma':=\overline{\mu}'-\overline{\nu}',
    \qquad
    \eta:=\sigma-\sigma'.
\]
Since \(\mu',\nu'\) are exactly \(k\)-wise indistinguishable, $M_{n\to k}\sigma'=0$ and therefore $M_{n\to k}\eta=M_{n\to k}\sigma=\gamma\theta_{k,k}$.

We lower-bound the total variation mass of \(\eta\). Starting with the fact $1=\sum_{a=0}^k \phi_k^{(k)}(a)\theta_{k,k}(a)$ and using the
adjoint relationship~\eqref{eq:T-M-adjoint-measures}, gives
\[
\begin{aligned}
    \gamma
    &=
    \sum_{a=0}^k \phi_k^{(k)}(a)\,\gamma\theta_{k,k}(a)              \\
    &=
    \sum_{a=0}^k \phi_k^{(k)}(a)(M_{n\to k}\eta)(a)              \\
    &=
    \sum_{j=0}^n (T_{n,k}\phi_k^{(k)})(j)\eta(j)                 \\
    &=
    \lambda_{n,k,k}
    \sum_{j=0}^n \phi_k^{(n)}(j)\eta(j).
\end{aligned}
\]
Thus
\[
    \left|
        \sum_{j=0}^n \phi_k^{(n)}(j)\eta(j)
    \right|
    =
    \frac{\gamma}{\lambda_{n,k,k}}.
\]
Using $\|\phi_k^{(n)}\|_\infty\le \sqrt{n+1}$,
we obtain
\[
    \frac{\gamma}{\lambda_{n,k,k}}
    \le
    \|\phi_k^{(n)}\|_\infty\,|\eta|(\{0,\dots,n\})
    \le
    \sqrt{n+1}\,|\eta|(\{0,\dots,n\}).
\]
Therefore
\begin{equation}
\label{eq:eta-mass-lower-bound}
    |\eta|(\{0,\dots,n\})
    \ge
    \frac{\gamma}{\lambda_{n,k,k}\sqrt{n+1}}.
\end{equation}
Since all distributions are symmetric, statistical distance between the
full distributions equals statistical distance between their
Hamming-weight laws. Moreover,
\[
    \eta
    =
    (\overline{\mu}-\overline{\mu}')
    -
    (\overline{\nu}-\overline{\nu}').
\]
Hence
\[
\begin{aligned}
    |\eta|(\{0,\dots,n\})
    &\le
    |\overline{\mu}-\overline{\mu}'|(\{0,\dots,n\})
    +
    |\overline{\nu}-\overline{\nu}'|(\{0,\dots,n\})              \\
    &=
    2\Delta(\mu,\mu')
    +
    2\Delta(\nu,\nu')                                           \\
    &\le
    4\max\{\Delta(\mu,\mu'),\Delta(\nu,\nu')\}.
\end{aligned}
\]
Combining this with~\eqref{eq:eta-mass-lower-bound}, we get
\[
    \max\{\Delta(\mu,\mu'),\Delta(\nu,\nu')\}
    \ge
    \frac{\gamma}{4\lambda_{n,k,k}\sqrt{n+1}}
    =
    \frac{\delta}{2\lambda_{n,k,k}\sqrt{n+1}L_k}.
\]
Finally, since
\[
    L_k
    =
    |\theta_{k,k}|(\{0,\dots,k\})
    =
    \frac{1}{k+1}\sum_{a=0}^k |\phi_k^{(k)}(a)|
    \le
    \left(
        \frac{1}{k+1}\sum_{a=0}^k (\phi_k^{(k)}(a))^2
    \right)^{1/2}
    =
    1,
\]
we conclude that
\[
    d
    \ge
    \frac{\delta}{2\lambda_{n,k,k}\sqrt{n+1}}.
\]
This holds for every symmetric exactly \(k\)-wise indistinguishable pair
\(\mu',\nu'\). Taking the infimum over all such pairs gives
\[
    \operatorname{Rep}_{n,k}(\mu,\nu)
    \ge
    \frac{\delta}{2\lambda_{n,k,k}\sqrt{n+1}}.
\]
\end{proof}

\end{document}
